\chapter{Dữ liệu và phương pháp}
\label{ch:phuongphap}

\section{Bộ dữ liệu pkdarabi/cardetection}

Bộ dữ liệu gồm 15 lớp: hai lớp đèn tín hiệu (Green Light, Red Light), mười hai lớp biển giới hạn
tốc độ và lớp Stop. Tổng cộng 4.969 ảnh và 6.012 đối tượng hợp lệ, chia sẵn thành ba tập như
Bảng~\ref{tab:split}.

\begin{table}[H]
\centering
\caption{Thống kê theo split}
\label{tab:split}
\small
\begin{tabular}{lrrrr}
\toprule
\textbf{Split} & \textbf{Số ảnh} & \textbf{Số đối tượng} & \textbf{Đối tượng/ảnh (TB)}
& \textbf{Diện tích bbox (px, median)} \\
\midrule
train & 3.530 & 4.298 & 1,22 & 8.802 \\
valid &   801 &   944 & 1,18 & 22.928 \\
test  &   638 &   770 & 1,21 & 8.532 \\
\bottomrule
\end{tabular}
\end{table}

Chất lượng nhãn rất sạch: 0 file nhãn thiếu, 0 đối tượng không hợp lệ, 0 dòng nhãn sai định dạng.
Lớp phổ biến nhất là Red Light (787 đối tượng), hiếm nhất là Speed Limit 10 (22 đối tượng).

\subsection{Hạn chế về tính đại diện}

Dữ liệu là ảnh cắt kích thước chuẩn hoá (median $416\times416$), thu từ nhiều nguồn video/camera
(có cả FisheyeCamera). Kết luận của đề tài giới hạn ở phạm vi bộ dữ liệu này và ở \emph{hành vi
của mô hình đã huấn luyện}, không khái quát cho mọi điều kiện đường thực tế.

\section{Công cụ đo --- YOLOv8s huấn luyện hướng-EDA}

Notebook nguồn suy ra cấu hình huấn luyện từ kết quả EDA rồi huấn luyện hai giai đoạn trên GPU
Tesla T4 (Ultralytics 8.4.95):

\begin{itemize}
  \item Mô hình nền \texttt{yolov8s.pt} (11,13 triệu tham số, 28,5 GFLOPs); dự phòng
  \texttt{yolov8n.pt} khi thiếu bộ nhớ.
  \item Giai đoạn 1: \texttt{imgsz=640}, batch 16, tối đa 60 epoch, \texttt{lr0=0,001}.
  \item Giai đoạn 2: \texttt{imgsz=768}, batch 10, tối đa 35 epoch, \texttt{lr0=0,0003}.
  \item Tối ưu AdamW + cosine LR; \textbf{không} lật ngang/dọc vì lật làm sai ngữ nghĩa biển báo
  có hướng.
\end{itemize}

Trọng số đã huấn luyện (\texttt{best.pt}) được đính kèm trong thư mục \texttt{models/} của đề
tài, phục vụ ứng dụng demo và việc kiểm chứng.

\section{Phương pháp phân tích thứ cấp}

Quy trình gồm ba bước, đều tự động và tái lập:

\begin{enumerate}
  \item \textbf{Trích xuất} (\texttt{src/extract\_notebook\_results.py}): đọc lại toàn bộ con số
  đã in trong notebook, cố định thành các bảng CSV chuẩn trong \texttt{results/tables/}. Không có
  giá trị nào nhập tay.
  \item \textbf{Kiểm định} (\texttt{src/audit\_analysis.py}): tính phơi nhiễm rò rỉ, khoảng cách
  đèn--biển (Mann--Whitney), và tương quan độ hiếm--hiệu năng (Spearman).
  \item \textbf{Trực quan hoá} (\texttt{src/make\_figures.py}): sinh hình từ bảng đã trích.
\end{enumerate}

\section{Chỉ số và kiểm định}

Hiệu năng đo bằng Precision, Recall, mAP@0,5 và mAP@0,5:0,95 (theo Ultralytics). Khoảng cách giữa
hai nhóm lớp kiểm định bằng Mann--Whitney U một phía; liên hệ độ hiếm--hiệu năng đo bằng tương
quan Spearman. Mọi bước cố định hạt giống \texttt{seed=42}.
